\section{Discussion}
\label{sec:discussion}

The central methodological finding of this work is the identification and formal characterization of the \textbf{Retrieval--Fusion Calibration Hazard}: a previously undiagnosed evaluation artifact that affects both decision-level classifier fusion and retrieval-augmented generation systems. The Hazard manifests through a common mechanism: when heterogeneous score distributions from different modalities or retrieval streams are combined via information-theoretic or heuristic rank fusion rules without prior probability calibration, the resulting fused scores inherit distortions that inflate apparent performance metrics.

\subsection{Stage I: Decision-Level Fusion Hazard Mechanism}

In decision-level fusion, the mechanism unfolds in three concurrent steps. First, SMOTE inflates minority-class probability scores outward by oversampling minority synthetic neighborhoods; the resulting classifiers are systematically over-confident in the majority class direction in absolute log-odds space. Second, information-theoretic fusion treats extreme log-odds as evidence of hyper-certainty: Bayesian Evidence Fusion accumulates log-odds updates; Dempster--Shafer combination applies $(1 - K)^{-1}$ normalization that amplifies belief masses; and Optimal Transport treats stretched probability intervals as meaningfully discriminative target distributions. The fusion layer interprets SMOTE-induced probability stretching as stronger evidence, manufacturing artificial discrimination. Third, Platt cross-calibration rescales probability outputs without altering their within-fold rank order: Platt's sigmoid preserves rank per fold ($\Delta \text{AUROC} = 0.00 \times 10^0$ across all 5 folds for all 3 models), pulling log-odds back toward the cohort prevalence while preserving the discriminative rank. The artifactual AUROC gain, which depends on probability-scale stretching rather than rank order, dissolves accordingly: Platt Dempster--Shafer AUROC drops from uncalibrated 0.7449 to 0.6831; Platt Optimal Transport AUROC drops from uncalibrated 0.7505 to 0.6636; and Platt Bayesian Evidence Fusion drops from uncalibrated 0.7438 to 0.6595.

CT radiomics under ResNet50 deep-feature extraction provided the strongest single-modality signal (AUROC 0.7037, Platt BSS +0.0515, AUPRC 0.3701). The dominance of M3 over M1 and M2 reflects three reinforcing mechanisms: imaging captures tumor heterogeneity directly through deep features corresponding to irregular tumor margins and necrotic core regions, which mechanistically upstream-drive metastatic propensity \cite{chi2026multimodal, nazari2020ct}; cross-sectional imaging is operationally available pre-operatively as part of routine staging, providing aligned features without patient-selection bias affecting genomic assays; and CT radiomics is not subject to SMOTE amplification due to the absence of a strong class-imbalance constraint in the imaging feature space. Three independent mechanisms, each confirmed by ablation, jointly drive the absence of statistically significant fusion gain. First, inter-modality residual error correlation with $\rho \in [0.3842, 0.4619]$ ($p < 0.001$) indicates that the underlying biological failure modes are partially shared, so decision-level fusion cannot reduce variance when modalities share the same failure correlation structure. Second, high Dempster--Shafer evidence conflict with $\bar{K} = 0.4381$ and 54.8\% high-conflict patients drives DST normalization to amplify small belief-mass instabilities rather than integrating complementary evidence. Third, SMOTE-induced probability distortion contributes a $+0.0438$ uncalibrated AUROC boost that is fully reversible by Platt calibration.

\subsection{Stage II: Retrieval--Fusion Calibration Hazard in RAG}

The Stage~II RAG experiments provide empirical proof that the Calibration--Fusion Hazard is not confined to decision-level classifier ensembles but extends directly to retrieval-augmented generation systems that fuse heterogeneous retrieval streams. When a medical RAG system combines BM25 (unbounded, right-skewed $[0, 15+]$), dense PubMedBERT cosine similarity (bounded $[-1,1]$), and radiomic feature similarity (bounded $[0,1]$) through ad-hoc normalization or Reciprocal Rank Fusion, the resulting fused relevance scores inherit severe probability distortion.

The empirical evidence reveals two key insights:
\begin{enumerate}
    \item \textbf{Catastrophic Miscalibration under Uncalibrated Rank Fusion:} Uncalibrated RRF produces severe overconfidence ($\text{ECE} = 0.5718, \text{Brier} = 0.5625$), rendering raw retrieval scores uninformative about true relevance. In contrast, Naive dense-only retrieval maintains well-calibrated confidence ($\text{ECE} = 0.0633, \text{Brier} = 0.2388$) because its score space is homogeneous and uncorrupted by heuristic rank mixing.
    \item \textbf{Partial Remediation via Calibration and Genuine Multimodal Gains:} Applying 5-fold out-of-fold Platt calibration reduces ECE to $0.3735$ (and Brier score to $0.3741$). While this represents a meaningful reduction in probability distortion, $\text{ECE} = 0.3735$ remains suboptimal, demonstrating that post-hoc monotonic scaling cannot fully eliminate the structural miscalibration introduced by multi-source rank fusion. Crucially, the genuine discriminative benefit of multimodal retrieval survives calibration: Platt-calibrated Multimodal RAG achieves a statistically significant RAGAS improvement over Naive RAG ($\Delta = +0.0481$, $p = 0.001$), with balanced class-stratified accuracy ($M_1$ Sensitivity = 0.7778, $M_0$ Specificity = 0.7685).
\end{enumerate}

\subsection{Agentic RAG Tradeoffs and Reasoning Dynamics}

The Agentic RAG architecture, which employs real LLM-driven ReAct tool-calling with structured Thought $\to$ Action $\to$ Observation reasoning traces, achieves comparable clinical correctness (0.7533; $M_1$ sensitivity 0.7222) but lower composite generation quality (RAGAS = 0.4870 vs. 0.5688 for calibrated multimodal hybrid RAG) due to lower Clinical Relevance (0.5393 vs. 0.6773). 

This empirical finding aligns with emerging literature on tool-calling dynamics in domain-specific medical LLMs \cite{singh2025agentic}: multi-round reasoning loops introduce extraneous meta-cognitive text that dilutes answer conciseness and term density on focused guideline corpora, without yielding commensurate accuracy gains over direct multimodal retrieval. In small, highly curated guideline settings, direct hybrid retrieval with calibrated evidence fusion offers a superior accuracy-to-efficiency profile compared to iterative agentic loops.

\subsection{Comparative Analysis with Existing ccRCC Benchmarks}

To contextualize the methodological advance of RenoFusion within the broader multimodal machine learning literature in renal oncology, Table~\ref{tab:comparative_analysis} provides a systematic comparative synthesis of our framework alongside seven landmark ccRCC predictive benchmarks published between 2021 and 2026 \cite{schulz2021multimodal, mahootiha2023multimodal, yang2024multimodal, ma2025genomic, zang2025multimodal, chi2026multimodal, xiao2026urofusion}.

\begin{table*}[t]
\centering
\caption{Comparative Analysis of RenoFusion Against Landmark ccRCC Multimodal Benchmarks.}
\label{tab:comparative_analysis}
\small
\renewcommand{\arraystretch}{1.2}
\setlength{\tabcolsep}{5pt}
\begin{tabularx}{\textwidth}{@{} >{\raggedright\arraybackslash}p{0.15\textwidth} c >{\raggedright\arraybackslash}p{0.17\textwidth} c >{\centering\arraybackslash}p{0.13\textwidth} L @{}}
\toprule
\textbf{Study / Benchmark} & \textbf{Year} & \textbf{Modalities Integrated} & \textbf{Cohort ($n$)} & \textbf{Reported Metric} & \textbf{Out-of-Fold Calibration \& Methodological Rigor} \\
\midrule
Schulz et al. \cite{schulz2021multimodal} & 2021 & Histology + CT + WES & 230 & C-index 0.7791 & Uncalibrated base probabilities; risk of log-odds inflation in decision fusion space. \\ \addlinespace[2pt]
Mahootiha et al. \cite{mahootiha2023multimodal} & 2023 & 3D CT + Clinical & 150 & AUROC 0.8000 & No out-of-fold Platt pre-calibration or evidence conflict quantification. \\ \addlinespace[2pt]
Yang et al. \cite{yang2024multimodal} & 2024 & Clinical + CT + US & 251 & AUROC 0.8200 & Imbalance-induced probability distortion unaddressed in downstream decision rules. \\ \addlinespace[2pt]
Ma et al. \cite{ma2025genomic} & 2025 & Multiphase CT + WSI & 274 & AUROC 0.8363 & Pre-fusion probability scaling omitted; unverified rank vs log-odds gain. \\ \addlinespace[2pt]
Zang et al. \cite{zang2025multimodal} & 2025 & Clinical + CT + WSI & 1,648 & C-index 0.8860 & Multi-center cohort, but calibration--fusion hazard left unexamined. \\ \addlinespace[2pt]
Chi et al. \cite{chi2026multimodal} & 2026 & CT Body Comp. + Omics & 180 & AUROC 0.7900 & Single-modality leaning; base-model probability miscalibration uncorrected. \\ \addlinespace[2pt]
Xiao et al. \cite{xiao2026urofusion} & 2026 & 3D CT + Pathology + Omics & Multi-center & AUROC 0.8400 & Product-of-experts architecture without mandatory Platt pre-calibration. \\ \addlinespace[2pt]
\textbf{RenoFusion (Ours)} & \textbf{2026} & \textbf{Clinical + Genomic + CT + LLM RAG} & \textbf{126 Aligned + 150 RAG} & \textbf{AUROC 0.7037 / RAGAS 0.5688} & \textbf{First to mandate Platt pre-calibration across both decision fusion and RAG; proves gains are artifactual ($p=0.442$); zero-leakage SMOTE; tri-factor decomposition; LLM-powered agentic RAG validation.} \\
\bottomrule
\end{tabularx}
\end{table*}

Prior state-of-the-art studies \cite{schulz2021multimodal, yang2024multimodal, zang2025multimodal, xiao2026urofusion} have achieved impressive discriminative scores by reporting raw uncalibrated fusion performance as headline metrics. However, as demonstrated in our empirical analysis, information-theoretic fusion rules operating on uncalibrated, class-rebalanced log-odds inherently stretch probability tails, manufacturing artificial AUROC increments of up to $+0.0468$ (e.g., uncalibrated Optimal Transport AUROC 0.7505 vs. Platt-calibrated 0.6636). 

RenoFusion advances beyond existing literature in five key dimensions:
\begin{enumerate}
    \item \textbf{First Calibration-Aware Evaluation Protocol:} RenoFusion is the first multimodal framework in renal oncology to enforce mandatory per-fold out-of-fold Platt scaling prior to both decision fusion and retrieval-score fusion, guaranteeing that reported metrics reflect genuine clinical discrimination rather than probability-scale distortion.
    \item \textbf{First LLM-Powered RAG Calibration Hazard Diagnosis:} RenoFusion provides the first empirical evidence that the Calibration--Fusion Hazard extends to retrieval-augmented generation, demonstrating catastrophic miscalibration ($\text{ECE} = 0.5718$) in uncalibrated rank fusion of heterogeneous retrieval streams.
    \item \textbf{Zero-Leakage Rebalancing Guardrails:} Unlike prior pipelines that apply SMOTE globally or across validation splits, RenoFusion strictly isolates synthetic oversampling to training folds, preventing probability tail leakage into validation sets.
    \item \textbf{Mechanistic Tri-Factor Decomposition:} While previous benchmarks present black-box fusion metrics, RenoFusion mathematically isolates why fusion gains dissolve, pinpointing inter-modality error correlation ($\rho \in [0.3842, 0.4619]$), high evidence conflict ($\bar{K} = 0.4381$), and SMOTE probability stretching ($+0.0438$ boost).
    \item \textbf{Real LLM Agentic RAG Validation:} Using Qwen2.5-3B-Instruct with ReAct tool-calling, RenoFusion validates that calibrated multimodal retrieval provides statistically significant improvement ($p = 0.001$) while agentic reasoning achieves balanced clinical correctness in domain-grounded clinical decision support.
\end{enumerate}

The methodological warning extends beyond ccRCC. Maier-Hein et al. \cite{maierhein2024metrics} cautioned that common calibration metrics can be highly model-dependent. Chen et al. \cite{chen2016calibrating} emphasized that learned risk scores must be subjected to post-hoc calibration regardless of the model's training-time probabilistic structure. Andersen et al. \cite{andersen2024impact} demonstrated that SMOTE systematically degrades calibration across 10 clinical datasets despite preserving rank-based discrimination. The present study unites these observations under a single empirical focus: within a single multimodal ccRCC pipeline spanning both classifier fusion and LLM-powered RAG, the joint application of class rebalancing, information-theoretic fusion, and ad-hoc retrieval-score normalization can manufacture predictive claims that disappear under proper cross-calibration. Per-modality out-of-fold probability calibration before fusion is a mandatory precondition for valid performance reporting in both decision-level and retrieval-augmented clinical AI systems.
